\clearpage
\section{Hooke’s Law from Free Energy}
\subsection{deriving Hookes law}
The differential of F is 
\begin{align}
    \dd{F}&=\dd{\ins{\frac{\lambda}{2}\epsilon_{ii}^2+\mu\epsilon_{ij}^2}}\\
    &= \lambda\epsilon_{kk}\dd{\epsilon_{ll}}+2\mu\epsilon_{ij}\dd{\epsilon_{ij}}\\
    &= \lambda\epsilon_{kk}\delta_{ij}\dd{\epsilon_{ij}}+2\mu\epsilon_{ij}\dd{\epsilon_{ij}}\\
    &= \ins{\lambda\epsilon_{kk}\delta_{ij}+2\mu\epsilon_{ij}}\dd{\epsilon_{ij}}
\end{align}
but by definition
\begin{equation}
    \dd{F}=\sigma_{ij}\dd{\epsilon_{ij}}
\end{equation}
so 
\begin{equation}
    \sigma_{ij}=\lambda\epsilon_{kk}\delta_{ij}+2\mu\epsilon_{ij}
\end{equation}
\subsection{comparison}
in class we saw 
\begin{equation}
    \sigma_{ij}=\frac{E}{1+\nu}\epsilon_{ij}+\frac{E\nu}{\ins{1+\nu}\ins{1-2\nu}}\epsilon_{kk}\delta_{ij}
\end{equation}
so we see the coefficients are 
\begin{equation}
    \mu = \frac{E}{2\ins{1+\nu}},\qquad \lambda=\frac{E\nu}{\ins{1+\nu}\ins{1-2\nu}}
\end{equation}